\section*{Chapter \ref{ch:dv:build-a-pricing-library} --- Build: A Pricing Library}

\begin{solution}{exo:dv:build-a-pricing-library:1}
Python salts \texttt{hash()} of strings per process, so the same instrument would get different seeds in two runs or two worker processes, and a bump
computed across processes would compare different random numbers. A checksum (CRC-32) of the identifier is the same everywhere and forever.
\end{solution}

\begin{solution}{exo:dv:build-a-pricing-library:2}
Each is priced at the volatility of its own strike, 100 at one year, which is a node. Bumping the 90 node changes the surface only between 80 and 100 (linear
interpolation in log-strike), and not at 100 itself. The knock-out's barrier at 80 would matter only to an engine that reads volatilities along the path.
\end{solution}

\begin{solution}{exo:dv:build-a-pricing-library:3}
In--out parity (the two barriers together are the vanilla); the digital call and put together pay one for sure, worth the discount factor; with no leverage and an
unreachable target the TARF is a strip of long puts, one per fixing.
\end{solution}

\begin{solution}{exo:dv:build-a-pricing-library:4}
With common random numbers the bumped and base runs differ only through the bump, so what remains is the Monte Carlo error of the sensitivity estimator itself,
which falls as $1/\sqrt n$. Reduce it with more paths, a control variate, a smoother payoff near the barriers, or a pathwise or likelihood-ratio \omterm{def:dv:greeks-and-the-hedging-pnl:greeks}{vega} (chapter 23)
behind the engine's own \texttt{greeks}.
\end{solution}

\begin{solution}{exo:dv:build-a-pricing-library:5}
Between expiries the surface interpolates total variance, $\sigma^2t$, not volatility, and within an expiry it interpolates in log-strike: a parallel shift of
every volatility is not exactly the sum of its node shifts. Central differences also leave second-order terms, which differ between one large bump and many
small ones.
\end{solution}

\begin{solution}{exo:dv:build-a-pricing-library:6}
The engine's price no longer follows the snapshot: bumps and scenarios do not reach it, its \omterm{def:dv:greeks-and-the-hedging-pnl:greeks}{Greeks} come out zero, and two threads pricing different snapshots
interfere. Parallel batches become wrong without any error being raised.
\end{solution}

\begin{solution}{exo:dv:build-a-pricing-library:7}
With 400\,000 paths the standard deviation over twenty seeds is 25.1 (mean 2\,231), against 50.8 with 100\,000: half, as $1/\sqrt n$ predicts for four times the paths.
\end{solution}

\begin{solution}{exo:dv:build-a-pricing-library:8}
Agreement with the old engine shows the two share their results, not that either is right: they may share a bug, a convention or a market input. The test must be
against an independent reference (a closed form, a converged grid) and an identity, over the whole range of inputs, including near barriers, expiries and
deep in and out of the money.
\end{solution}

\begin{solution}{pb:dv:build-a-pricing-library:1}
\textbf{1.} Call: Analytic; American put: Tree; knock-out: ClosedForm (daily monitoring by the Broadie--Glasserman--Kou shift); \omterm{def:dv:variance-swaps-and-volatility-derivatives:vs}{variance swap}: ClosedForm
(replication on the surface); \omterm{def:dv:autocallables:autocallable}{autocallable}: MonteCarlo. The registry returns the first registered engine that supports the pair.
\textbf{2.} $-928\,523$.
\textbf{3.} Analytic and tree to $2\times10^{-6}$; the grid 0.0024 (200 nodes); Monte Carlo within its standard error.
\textbf{4.} Read the volatility from the snapshot's surface at its own strikes and dates, never from a stored number.
\textbf{5.} It raises a pricing error naming the instrument and the model.
\textbf{6.} Call 0, American put $-88$, knock-out 0, \omterm{def:dv:variance-swaps-and-volatility-derivatives:vs}{variance swap} $-90$, \omterm{def:dv:autocallables:autocallable}{autocallable} $+2\,249$: 2\,072.
\textbf{7.} Its engine reads the total variance at 90, the midpoint of the knock-in and the trigger, and the note is worth less to its holder when volatility
rises; the firm is short the note, so long its \omterm{def:dv:greeks-and-the-hedging-pnl:greeks}{vega}.
\textbf{8.} Linear interpolation in log-strike between 90 and 100 puts about half of a 95 strike's weight on the 90 node.
\textbf{9.} $-90$: the strip of options that replicates variance spans all strikes, so every node carries part of its \omterm{def:dv:greeks-and-the-hedging-pnl:greeks}{vega}.
\textbf{10.} 2\,275 at one year, 1\,464 at two and 3\,008 at three per point, and nothing before one year.
\textbf{11.} 51 with common random numbers, 217 without.
\textbf{12.} The bumped and base prices share their errors, which cancel in the difference; without sharing, the errors of two independent runs add and are divided
by the bump.
\textbf{13.} About half: 25.
\textbf{14.} So that the price of one trade does not change when other trades are added to or removed from the book, and so that Book 6 can reprice any subset.
\textbf{15.} The book is long volatility and long a crash ($+161.8$ thousand for $-20\%$ with ten points more volatility), and loses on rallies ($-24.0$ thousand at
$+10\%$) and on falling volatility ($-34.9$ thousand).
\textbf{16.} To about a hundred: the Monte Carlo part alone has a standard deviation of 51.
\textbf{17.} Reference prices and cross-engine checks run nightly, a record of the engine and seed used for each trade, the snapshot stored with each result, and
limits on unexplained changes.
\textbf{18.} A property test that bumps the snapshot and requires the price to change for every instrument whose value depends on the bumped factor.
\textbf{19.} 2\,072 per volatility point (\omterm{def:dv:autocallables:autocallable}{autocallable} $+2\,249$, American put $-88$, \omterm{def:dv:variance-swaps-and-volatility-derivatives:vs}{variance swap} $-90$), with a standard deviation of 51 from Monte Carlo
noise with common random numbers and 217 without.
\textbf{20.} Pricing as a pure function of an immutable snapshot, so that a bumped snapshot reprices anything, with shared random numbers for simulation.
\end{solution}

\begin{solution}{iq:dv:build-a-pricing-library:1}
So that each varies independently: a new product needs its terms and one engine; a new method prices every product it supports; two engines check each other; the
model's assumptions are explicit and swappable.
\iqlookfor{independent variation, cross-checking.}
\end{solution}

\begin{solution}{iq:dv:build-a-pricing-library:2}
Common random numbers: the same seed and grid for every repricing of one instrument, so that bumps difference out the simulation error. Then smoother payoffs and
direct \omterm{def:dv:greeks-and-the-hedging-pnl:greeks}{Greeks}.
\iqlookfor{the same seed per instrument.}
\end{solution}

\begin{solution}{iq:dv:build-a-pricing-library:3}
Immutable snapshots for risk: pricing is a pure function, parallel batches are safe, scenarios are new snapshots, results are reproducible from stored inputs.
Observers and caching suit a trading screen that reprices as quotes tick; they complicate thousands of shifted markets and parallel evaluation.
\iqlookfor{purity and reproducibility against incremental recalculation.}
\end{solution}

\begin{solution}{iq:dv:build-a-pricing-library:4}
Against closed forms and independent \omterm{def:dv:build-a-pricing-library:reference}{reference pricers}, against other engines, by identities (parity, in--out), by limits (zero volatility, far strikes), by
convergence in its own parameters (steps, paths), and by regression tests pinned to certified numbers.
\iqlookfor{independent references and convergence.}
\end{solution}

\begin{solution}{iq:dv:build-a-pricing-library:5}
Put every engine on the same surface object with node risk factors, bump each node in turn on a copy of the snapshot, reprice every position with its own model
and engine, with common random numbers for simulations, and sum by node. Check that the nodes sum to the parallel \omterm{def:dv:greeks-and-the-hedging-pnl:greeks}{vega}.
\iqlookfor{one surface, node bumps, common random numbers.}
\end{solution}

\begin{solution}{iq:dv:build-a-pricing-library:6}
Freeze reference values from the Python library for a set of cases; port the same arithmetic, operation by operation (grids, boundary values, solve order); test
to $10^{-9}$ in both languages in continuous integration; only then optimise, keeping the tests.
\iqlookfor{reference values and operation-by-operation equivalence.}
\end{solution}
